% introduction_generative.tex
\subsection{Generative Language Models as Respondents, Raters, and Test Developers}

Others prompt generative models to act as respondents. \citet{liu-etal-2025} calibrated six models' college-algebra answers with item response theory: the best model's difficulty parameters tracked human calibrations (Spearman $\rho = .87$, GPT-3.5); mixing 50 humans with as many resampled synthetic respondents raised difficulty-ordering recovery from $\rho = .89$ to .93 at higher root-mean-square error; and far narrower ability distributions (standard deviations 0.29--0.58 versus the human 0.98) rule out any single model as a current substitute. The ``virtual respondents'' of \citet{lim-etal-2025}---mediating characteristics letting one trait yield different item answers---screened machine-generated Big Five, Schwartz values, and character-strength items into sets in the top 1\% (Big Five) to 13\% (Schwartz and VIA) of possible selections on human-derived convergent-validity criteria without claiming to replicate human psychology; pre-data vetting is feasible but only as faithful as its simulation.

Educational measurement instead predicts item parameters from item text: a 37-article (46-study) systematic review of large-scale assessment traces difficulty modeling from hand-crafted linguistic features to transformers, the strongest correlating .87 with empirical difficulty at root-mean-square errors as low as 0.165 \citep{peters-etal-2025}. The signal is uneven: fine-tuned transformers predicted 1{,}119 operational English Language Arts items' banked difficulty at correlations near .72, yet on discrimination a gradient-boosted tree over structured item attributes ($r = .54$) beat every fine-tuned model \citep{han-etal-2025}; a generative model fine-tuned to reproduce response probabilities at twenty discrete ability levels, however, recovers both parameters through synthetic item characteristic curves, discrimination especially well \citep{ormerod-2026}.

Generative models have also been validated as raters of constructs in text: in two stages---iterative prompt development, then confirmatory testing on manually coded held-out data---\citet{bunt-etal-2025} found GPT-4o excellent on directly observable phenomena (reported speech and conversational repair, $F_{1} = .91$ for the best prompts on both) but only moderate on an inferential ordinal construct, harm in healthcare complaints (weighted $\kappa = .49$--$.57$); valid machine coding requires actively establishing semantic, predictive, and content validity. \citet{eberhardt-etal-2025} applied classical scale construction---theory-driven item pool, automated selection, confirmatory factor analysis, reliability estimation---to ratings from a locally hosted model rather than human judges: scored over 1{,}131 automatically transcribed therapy sessions, their eight-item ``LLM rating scale'' for patient engagement showed excellent internal consistency ($\omega = .95$) and theoretically coherent correlations with patient motivation ($r = .41$), therapist-rated between-session effort ($r = .39$), and later symptom outcome ($r = -.30$)---direct precedent that psychometric machinery transfers to language-model output.

A practical guide spans whole-instrument construction: item development (training-data quality, multi-model comparison, output validation) and ``data-less'' calibration, where sentence encoders check construct boundaries and item--definition alignment, embedding-similarity matrices are factor-analyzed, and, absent any sample, fit rests on model-free residual diagnostics (root-mean-square residual, common part accounted for) \citep{suarez-alvarez-etal-2026}; yet it cautions that embedding-based loadings---mostly positive cosine similarities---do not yet differentiate reverse-keyed items as conventional loadings do, leaving the human-response factor structure the gold standard for semantic structures. With a structured zero-shot prompt, \citet{grobelny-etal-2025} had GPT-4 adapt an organizational-commitment scale from English to Polish; among 785 working adults randomized to either version, the adaptation equaled or surpassed the established human one in factor structure, measurement invariance, and convergent and discriminant validity, its internal consistency statistically lower ($\alpha = .89$ versus $.92$) yet within accepted standards.

These routes---simulated respondents \citep{liu-etal-2025,lim-etal-2025}, parameter prediction from text \citep{peters-etal-2025}, machine coding and rating scales \citep{bunt-etal-2025,eberhardt-etal-2025}, and AI-assisted construction and adaptation \citep{suarez-alvarez-etal-2026,grobelny-etal-2025}---promise much psychometric work before any respondent is recruited, yet each inherits its model's run-to-run stochasticity, prompt sensitivity, and black-box opacity \citep{bunt-etal-2025}. Semantic factor analysis complements them all: deterministic and exactly reproducible given a fixed embedding matrix, it emulates no behavior but answers a question simulation does not---how a scale's language is organized.
